\section{Setup and Notation}\label{sec:setup}

This section introduces the setup, notation, and estimators used throughout the paper.
Let $Y_i$ denote the outcome variable, $X_i$ the running variable, and $T_i \defeq \ind\{X_i \geq c\}$ the treatment indicator under a sharp regression discontinuity design with cutoff $c$.
Observed outcomes relate to potential outcomes via $Y_i = Y_i(0) + [Y_i(1) - Y_i(0)] T_i$, where $Y_i(0)$ is the baseline outcome and $Y_i(1) - Y_i(0)$ the unit-level treatment effect.
In addition to the running variable, each unit carries a continuous predetermined covariate $Z_i \in \mathcal{Z} \subseteq \mathbb{R}$.

\begin{assumption}[Predetermined Covariate]\label{as:predetermined}
	The covariate is predetermined, i.e., realized before treatment assignment and invariant to it: $Z_i(1) = Z_i(0)$ for all $i$.
\end{assumption}

This paper studies covariate adjustment for a single continuous predetermined covariate $Z_i$ in the standard RD setup.
The key concern is that the conditional density $f_{Z|X}(z|x)$ may vary with the running variable, so that the composition of units along $Z$ changes as one moves away from the cutoff.
When this compositional variation interacts with heterogeneity in the regression function across $z$, pooling all observations introduces additional smoothing bias beyond what arises from the outcome regression alone.

A leading motivating example is geographic regression discontinuity designs, where $X_i$ measures distance to a treatment boundary and $Z_i$ denotes position along that boundary.
In this setting, the conditional density of $Z$ given $X$ can vary dramatically: urban centers generate sharp density peaks at specific locations along the boundary, while rural areas are sparsely populated.
As distance to the boundary changes, the relative weight of these dense and sparse regions shifts, creating steep compositional gradients.
However, the mechanism is not specific to geography --- any RD application in which a continuous covariate has a conditional density that varies with the running variable faces the same issue.

\subsection{Regression Functions and Target Parameter}\label{sec:setup:regression}

Define the long regression functions conditional on both the running variable and the covariate:
\begin{align}
	m(x,z) &\defeq \E\left[Y_i(0) \given X_i = x,\, Z_i = z\right],\\
	\tau(x,z) &\defeq \E\left[Y_i(1) - Y_i(0) \given X_i = x,\, Z_i = z\right].
\end{align}
These are the group-specific baseline outcome and treatment effect functions from the long regression model
\begin{equation}\label{eq:long_regression}
	Y_i = m(X_i, Z_i) + \tau(X_i, Z_i)\, T_i + e_i,
\end{equation}
with $\E[e_i \mid X_i, Z_i] = 0$.

The pooled (short) regression functions that average over the covariate are
\begin{align}
	\bar{m}(x) &\defeq \E\left[m(X_i, Z_i) \given X_i = x\right] = \int m(x,z)\, f_{Z|X}(z|x)\, dz,\\
	\bar{\tau}(x) &\defeq \E\left[\tau(X_i, Z_i) \given X_i = x\right] = \int \tau(x,z)\, f_{Z|X}(z|x)\, dz,
\end{align}
where $f_{Z|X}(z|x)$ is the conditional density of $Z_i$ given $X_i = x$.
The pooled regression model is then
\begin{equation}\label{eq:pooled_regression}
	Y_i = \bar{m}(X_i) + \bar{\tau}(X_i)\, T_i + u_i,
\end{equation}
where $u_i$ absorbs the covariate-driven heterogeneity.

The canonical RD parameter, the average treatment effect at the cutoff, can be written as a weighted average of covariate-specific treatment effects:
\begin{equation}\label{eq:tau_rd}
	\tau(c) \defeq \bar{\tau}(c) = \int \tau(c,z)\, f_{Z|X}(z|c)\, dz.
\end{equation}
More generally, this paper considers convex weighted averages of covariate-specific treatment effects along the boundary:
\begin{equation}\label{eq:tau_w}
	\tau_w \defeq \int \tau(c,z)\, w(z)\, dz,
\end{equation}
for a known, nonnegative weight function $w$ satisfying $\int w(z)\, dz = 1$.
Since $w$ is chosen by the researcher rather than estimated, it does not contribute to the estimation problem.
The canonical RD parameter corresponds to $w(z) = f_{Z|X}(z|c)$.

\paragraph{Fixed-design framework.}
For the minimax inference developed in subsequent sections, both $X_i$ and $Z_i$ are treated as part of the fixed design: inference is conducted conditional on the realized values $(X_1, Z_1), \ldots, (X_n, Z_n)$, following \textcite{armstrong_simple_2018}.
At the sample level, the weighted estimand takes the form $\tau_w = \sum_i \lambda_i \tau(c, Z_i)$ with pre-specified weights~$\lambda_i$, a linear functional of the bivariate regression function evaluated at the design points.
This sidesteps the conditional density $f_{Z|X}$ and its derivatives as nuisance objects in the inference problem; the population-level definitions above motivate the covariate adjustment problem, but the minimax machinery operates on the bivariate function $m(x,z)$ at the observed $(X_i, Z_i)$.
When $Z_i$ are viewed as draws from some distribution $F_Z$, the sample-level estimand approximates $\int \tau(c,z)\, w(z)\, dz$ as $n \to \infty$, but the finite-sample inference does not depend on this.


\subsection{Pooled Estimation --- Base RD}\label{sec:setup:base}

The standard approach to estimating $\tau(c)$ is the base RD estimator, which pools all observations regardless of the covariate.
For local constant estimation ($q = 0$), this reduces to comparing sample means within a bandwidth $h$ of the cutoff:
\begin{equation}\label{eq:base_rd}
	\hat{\tau}(h) = \frac{\sum_{i:\, X_i \geq c} K_h(X_i)\, Y_i}{\sum_{i:\, X_i \geq c} K_h(X_i)} - \frac{\sum_{i:\, X_i < c} K_h(X_i)\, Y_i}{\sum_{i:\, X_i < c} K_h(X_i)},
\end{equation}
where $K_h(X_i) \defeq K\!\left(\frac{X_i - c}{h}\right)$ for a kernel function $K$ with bounded support.
With a uniform kernel, this is simply the difference in sample means within the window $[c - h,\, c + h]$.

The base RD estimator implicitly targets the pooled regression function $\bar{m}(x) + \bar{\tau}(x) \ind\{x \geq c\}$.
Its smoothing bias therefore depends on the derivatives of $\bar{m}$ and $\bar{\tau}$, which inherit cross-terms from covariate heterogeneity and compositional changes, as discussed in \cref{sec:setup:pooling_bias}.


\subsection{Segmented Estimation --- Covariate-Adjusted RD}\label{sec:setup:segmented}

To avoid pooling across different values of the continuous covariate, I propose a segmented estimator.
The idea is to partition the covariate space $\mathcal{Z}$ into $J$ segments $\{S_1, \ldots, S_J\}$ and estimate the treatment effect separately within each segment using a local constant base RD estimator, then form a weighted average.

\paragraph{Segmentation.}
Let $\{S_j\}_{j=1}^J$ be a partition of $\mathcal{Z}$ into intervals of width $\delta$.
A natural choice, which I adopt throughout, is to set the segment width equal to the bandwidth: $\delta = h$.
This creates $h \times h$ ``squares'' in the $(X, Z)$-plane on each side of the cutoff.

\paragraph{Within-segment estimation.}
For each segment $S_j$, define the within-segment local constant RD estimator:
\begin{equation}\label{eq:tau_j}
	\hat{\tau}_j(h) = \frac{\sum_{i:\, X_i \geq c} K_h(X_i)\, \ind\{Z_i \in S_j\}\, Y_i}{\sum_{i:\, X_i \geq c} K_h(X_i)\, \ind\{Z_i \in S_j\}} - \frac{\sum_{i:\, X_i < c} K_h(X_i)\, \ind\{Z_i \in S_j\}\, Y_i}{\sum_{i:\, X_i < c} K_h(X_i)\, \ind\{Z_i \in S_j\}},
\end{equation}
which compares average outcomes on both sides of the cutoff among units in segment $S_j$.

\paragraph{Aggregation.}
The segmented RD estimator combines the within-segment estimates using chosen weights:
\begin{equation}\label{eq:tau_seg}
	\hat{\tau}_S(h) = \sum_{j=1}^{J} w_j\, \hat{\tau}_j(h),
\end{equation}
where the weights $w_j \geq 0$ satisfy $\sum_j w_j = 1$ and are chosen by the researcher to target the parameter $\tau_w$ in \cref{eq:tau_w}.

\begin{remark}[Relation to Fixed Effects Estimation]\label{rem:fe}
	For a fixed number of segments $J$, the segmented estimator \eqref{eq:tau_seg} is closely related to a fixed effects (FE) estimator that includes segment indicators.
	By the Frisch--Waugh--Lovell theorem \parencite[see also][for the extension to standard errors]{ding_frischwaughlovell_2021}, a regression of $Y_i$ on $T_i$ and segment dummies $\ind\{Z_i \in S_j\}$ within the bandwidth yields the same treatment effect estimate as the weighted average \eqref{eq:tau_seg} with specific weights.
	In particular, the OLS fixed effects estimator corresponds to weights that are proportional to the product of the within-segment treated and control sample sizes, i.e., inverse propensity score (variance-minimizing) weights \parencite{negi_revisiting_2021}.
	The segmented estimator \eqref{eq:tau_seg} generalizes this by allowing arbitrary convex weight choices.
\end{remark}


\subsection{Pooling Bias with Continuous Covariates}\label{sec:setup:pooling_bias}

To see how pooling across a continuous covariate induces bias, consider the derivatives of the pooled regression function.
Taking the first derivative of $\bar{m}(x) = \int m(x,z)\, f_{Z|X}(z|x)\, dz$ yields:
\begin{equation}\label{eq:mbar_deriv}
	\bar{m}'(x) = \int \frac{\partial m(x,z)}{\partial x}\, f_{Z|X}(z|x)\, dz + \int m(x,z)\, \frac{\partial f_{Z|X}(z|x)}{\partial x}\, dz.
\end{equation}
The first term reflects the average within-covariate-value slope of the regression function.
The second term arises from compositional changes: as $x$ varies, the conditional distribution of $Z$ shifts, reweighting the regression function.
Even if $m(x,z)$ is smooth in $x$ for every $z$, the pooled function $\bar{m}(x)$ can exhibit additional curvature due to changes in $f_{Z|X}(z|x)$.

For the local constant estimator, the smoothing bias depends on $\bar{m}'(x)$.
Steep or irregular conditional densities $f_{Z|X}(z|x)$ --- as arise in geographic settings with highly uneven population density --- inflate this derivative and thereby increase the pooling bias.
The segmented estimator mitigates this mechanism: within each segment $S_j$, the estimator targets the conditional mean $\E\left[Y_i(0) \given X_i = x,\, Z_i \in S_j\right]$, which averages over only a narrow range of $Z$-values.
This reduces the compositional term, since $m(x,z)$ varies by at most $M_z \cdot \delta$ within a segment of width $\delta$.
Unlike the discrete case, where stratification fully eliminates pooling bias, segments of positive width retain a residual compositional contribution that vanishes as $\delta \to 0$.
